% HodgeCY II Hilbert-Burch theorem/proof package.
% This file is a manuscript asset, not a standalone LaTeX document.

\begin{theorem}[Eight-plane line skeleton]\label{thm:hodgecy-ii-eight-plane-line-skeleton}
Let \(S=k[x_0,x_1,x_2,x_3]\), and let
\(L_1,\ldots,L_8\in S_1\) be distinct linear forms such that every
pair has rank \(2\) and every triple has rank \(3\).  Equivalently, no
three of the corresponding planes in \({\bf P}^3\) contain a common
line.  Put
\[
  F=\prod_{i=1}^8 L_i,
  \qquad
  I_C=(F/L_1,\ldots,F/L_8).
\]
Then
\[
  I_C=\bigcap_{1\le i<j\le 8}(L_i,L_j).
\]
Thus \(I_C\) is the saturated reduced ideal of the union \(C\) of the
\(\binom{8}{2}=28\) pairwise intersection lines.  Moreover \(I_C\) has
Hilbert--Burch resolution
\[
  0 \longrightarrow S(-8)^7
    \xrightarrow{\,M\,}
  S(-7)^8
    \longrightarrow I_C
    \longrightarrow 0,
\]
where
\[
M=
\begin{pmatrix}
L_1&0&0&0&0&0&0\\
0&L_2&0&0&0&0&0\\
0&0&L_3&0&0&0&0\\
0&0&0&L_4&0&0&0\\
0&0&0&0&L_5&0&0\\
0&0&0&0&0&L_6&0\\
0&0&0&0&0&0&L_7\\
-L_8&-L_8&-L_8&-L_8&-L_8&-L_8&-L_8
\end{pmatrix}.
\]
\end{theorem}

\begin{proof}
Write \(G_i=F/L_i\).  Each \(G_m\) belongs to every pair ideal
\((L_i,L_j)\): if \(m\notin\{i,j\}\), then \(G_m\) contains both
factors \(L_i,L_j\); if \(m=i\), it contains \(L_j\); and if \(m=j\),
it contains \(L_i\).  Hence
\[
  I_C\subseteq \bigcap_{i<j}(L_i,L_j).
\]

For the reverse inclusion and saturation, use the codimension-two
basic-double-link induction.  For \(s\ge 2\), let
\[
  J_s=\bigcap_{1\le i<j\le s}(L_i,L_j),
  \qquad
  P_s=\prod_{i=1}^s L_i.
\]
The case \(s=2\) is \(J_2=(L_1,L_2)\).  Assume
\[
  J_{s-1}=(P_{s-1}/L_1,\ldots,P_{s-1}/L_{s-1}).
\]
Since no three planes contain a line, the plane \(L_s=0\) contains no
line component of the old codimension-two skeleton.  The elementary
basic-double-link exact sequence, equivalently Proposition 2.6 of
Geramita--Harbourne--Migliore applied in this codimension-two
situation, gives
\[
  J_s=L_sJ_{s-1}+(P_{s-1}).
\]
Substituting the induction hypothesis yields
\[
  J_s=
  (L_sP_{s-1}/L_1,\ldots,L_sP_{s-1}/L_{s-1},P_{s-1})
  =(P_s/L_1,\ldots,P_s/L_s).
\]
Taking \(s=8\) proves the displayed ideal identity.  The right-hand
side is an intersection of distinct line primes, so it is saturated and
reduced.

It remains to record the resolution.  The maximal minors of the
displayed \(8\times 7\) matrix \(M\), up to the usual alternating
signs, are exactly the generators \(G_1,\ldots,G_8\).  The columns give
the syzygies
\[
  L_iG_i-L_8G_8=0,\qquad i=1,\ldots,7.
\]
The generators have no common factor, so the ideal has height \(2\).
The Hilbert--Burch theorem therefore gives the stated free resolution.
\end{proof}

\begin{corollary}[Hilbert series of the line skeleton]\label{cor:hodgecy-ii-line-skeleton-hilbert-series}
For the curve \(C\) of Theorem~\ref{thm:hodgecy-ii-eight-plane-line-skeleton},
\[
  \operatorname{Hilb}_{S/I_C}(t)
  =
  \frac{1-8t^7+7t^8}{(1-t)^4}.
\]
\end{corollary}

\begin{proof}
Take Hilbert series in the Hilbert--Burch resolution of
Theorem~\ref{thm:hodgecy-ii-eight-plane-line-skeleton}.
\end{proof}

\begin{theorem}[Regular quartic block section]\label{thm:hodgecy-ii-regular-quartic-block-section}
Let \(C\) be as in Theorem~\ref{thm:hodgecy-ii-eight-plane-line-skeleton}.
Let \(Q\in S_4\) be a quartic whose image in \(S/(L_i,L_j)\) is nonzero
for every associated line prime \((L_i,L_j)\) of \(S/I_C\).  Then \(Q\)
is a non-zero-divisor on \(S/I_C\).  If
\[
  I_B=I_C+(Q),
\]
then there is an exact sequence
\[
  0 \longrightarrow (S/I_C)(-4)
  \xrightarrow{\cdot Q}
  S/I_C
  \longrightarrow S/I_B
  \longrightarrow 0.
\]
For the frozen HodgeCY II arrangements \(84\) and \(84a\), the verifier
checks all \(28\) line restrictions and certifies that \(Q\) restricts
to a nonzero squarefree quartic on every associated line.
\end{theorem}

\begin{proof}
The ideal \(I_C\) is reduced and saturated with associated primes
exactly the line primes \((L_i,L_j)\).  A homogeneous element is a
non-zero-divisor on \(S/I_C\) precisely when it is contained in none of
these associated primes.  The stated nonvanishing of each line
restriction is exactly this condition.  Multiplication by \(Q\) is
therefore injective on \(S/I_C\), giving the displayed exact sequence.
The final assertion is the frozen HodgeCY II computation recorded in
the theorem evidence bundle.
\end{proof}

\begin{corollary}[Mapping-cone resolution]\label{cor:hodgecy-ii-block-mapping-cone}
For the block scheme \(B=C\cap V(Q)\) of
Theorem~\ref{thm:hodgecy-ii-regular-quartic-block-section}, \(S/I_B\)
has free resolution
\[
0
\longrightarrow S(-12)^7
\longrightarrow S(-11)^8 \oplus S(-8)^7
\longrightarrow S(-7)^8 \oplus S(-4)
\longrightarrow S
\longrightarrow S/I_B
\longrightarrow 0.
\]
\end{corollary}

\begin{proof}
Apply the mapping cone to multiplication by \(Q\) on the
Hilbert--Burch resolution of \(S/I_C\).  Since multiplication by \(Q\)
is injective on \(S/I_C\), the resulting complex is exact.
\end{proof}

\begin{corollary}[Structural Hilbert collapse]\label{cor:hodgecy-ii-structural-hilbert-collapse}
For the verified HodgeCY II block schemes attached to \(84\) and
\(84a\),
\[
  \operatorname{Hilb}_{S/I_B}(t)
  =
  \frac{(1-t^4)(1-8t^7+7t^8)}{(1-t)^4}.
\]
Consequently
\[
  H_B(d)=1,4,10,20,34,52,74,92,105,112,112,\ldots,
\]
with stabilization at degree \(9\).
\end{corollary}

\begin{proof}
Use the exact sequence of
Theorem~\ref{thm:hodgecy-ii-regular-quartic-block-section}.  Thus
\[
  \operatorname{Hilb}_{S/I_B}(t)
  =(1-t^4)\operatorname{Hilb}_{S/I_C}(t).
\]
Substitute Corollary~\ref{cor:hodgecy-ii-line-skeleton-hilbert-series}
and expand the coefficients.
\end{proof}

\begin{corollary}[Degree-eight cohomology and block-evaluation deficiency]\label{cor:hodgecy-ii-degree-eight-deficiency}
For the verified HodgeCY II block schemes attached to \(84\) and
\(84a\),
\[
  h^1({\bf P}^3,\mathcal I_B(8))=7.
\]
Equivalently, the degree-eight evaluation cokernel has dimension
\[
  \epsilon_B=7
\]
for the verified block scheme.
\end{corollary}

\begin{proof}
The vector space of degree-eight forms in \(S\) has dimension
\(\binom{11}{3}=165\).  Since \(B\) has degree \(112\) and
Corollary~\ref{cor:hodgecy-ii-structural-hilbert-collapse} gives
\(H_B(8)=105\), the standard exact sequence
\[
0\to H^0(\mathcal I_B(8))\to H^0(\mathcal O_{{\bf P}^3}(8))
\to H^0(\mathcal O_B)\to H^1(\mathcal I_B(8))\to 0
\]
gives
\[
  h^1(\mathcal I_B(8))=112-105=7.
\]
The evaluation map on degree-eight forms has image dimension
\(H_B(8)=105\), so its cokernel also has dimension \(7\).
This proves the block-scheme deficiency only.  It is not a promotion
of the classical nodal defect.
\end{proof}

\begin{proposition}[HodgeCY II non-determination witness]\label{prop:hodgecy-ii-nondetermination}
For the frozen witness pair \(84/84a\), the verified source data give
\[
  HA_{\bf Q}^{(2)}(84)\cong HA_{\bf Q}^{(2)}(84a),
  \qquad
  HA_{\bf Z}^{(2)}(84)\not\cong HA_{\bf Z}^{(2)}(84a),
\]
while the verified block schemes have identical structural
Hilbert/evaluation data through degree \(8\), including
\(\epsilon_B=7\).  Therefore the verified degree-eight block-evaluation
data do not determine the integral source-assembly type on the witness
pair \(84/84a\).
\end{proposition}

\begin{proof}
The source-lattice certificate records equal rational source rank
\(26\) and equal rational \(H_1\)-rank \(2\), but different integral
Smith forms
\[
  (1^{23},2,6,12)
  \quad\hbox{and}\quad
  (1^{21},2,4,4,4,12).
\]
Corollaries~\ref{cor:hodgecy-ii-structural-hilbert-collapse} and
\ref{cor:hodgecy-ii-degree-eight-deficiency} prove that the verified
block Hilbert/evaluation signatures coincide.  This is a
non-determination statement on this witness pair only.  It does not
assert a reverse implication, a canonical source-to-evaluation map, an
integral evaluation lattice, an LMHS identification, or a classical
defect promotion.
\end{proof}

